\scanpage{65}
\deleted{\unsure{équivaut à celle qu'on va écrire}}
\[
\deleted{\boxed{\underline\lambda^{\unsure{g|\sigma}}\lambda^{-1}=}}
\]

Il y a une solution triviale
\[
\lambda=\frac1L\qquad\text{i.e.}\quad L'(t)=1.
\]
Mais nous désirons imposer que $\lambda\deleted{\unsure{L}}$ soit de ``niveau'' \unsure{strict.} inf. à $L$.

Pour préciser ce point, nous supposerons
\deleted{\fbox{$L(0)=1$,} que les racines et pôles de $L$ sont en valeurs absolues de la forme
\[
p^{-\frac i2}\quad\text{avec}\quad \unsure{4|2}\rho\leq i\leq\unsure{4|2}\rho+2n\qquad\text{i.e.}\quad\rho\leq\frac i2\leq\rho+n,
\]
\deleted{et \unsure{modulo multiplication} par des puissances de $p^{1/2}$, des valeurs \uuline{propres inverses}}
(et on veut se restreindre aux $\lambda$ ayant \uline{la \unsure{$=$|même} propriété}, \deleted{et de plus} (de poids compris entre $\rho$ et $\rho+n$) et de plus de ``niveau'' \deleted{$\unsure{\rho+n+1}$}}

\margin{\unsure{et sont de la forme} $(p^{1/2})^{\unsure{\nu|n}}\xi_i$, où $\nu$ un entier, $\xi_i$ un \uline{entier alg.}}

\diagram{\begin{tikzcd}[column sep=small]
L(t)={} & {L_\rho(t)\ldots L_{\rho\deleted{\unsure{\delta|2}}+n}(t)} \arrow[loop below,distance=1.4em,"{\deleted{\unsure{=L}}}"'] & {L'_{\rho+n-1}(pt)L'_{\rho+n-2}(p^2t)\ldots L'_\rho(p^nt)}
\end{tikzcd}}
\deleted{\unsure{et de}}
où les valeurs propres inverses des $\unsure{L_{i/2}|L_i}$ sont des entiers alg en \unsure{val.} absolue $p^{i/2}$, et de \unsure{$=$|même} pour $L'_i$ (donc $L'_{\rho+n-i}(p^it)$ a des valeurs propres inverses qui sont des entiers alg de val. absolue $p^{\frac{\rho+n-i}{2}}p^i=p^{\frac{\rho+n+i}{2}}$.

\margin{64}
